
\documentclass{article}
\usepackage{colortbl}
\usepackage{makecell}
\usepackage{multirow}
\usepackage{supertabular}

\begin{document}

\newcounter{utterance}

\centering \large Interaction Transcript for game `dond', experiment `coop', episode 4 with Qwen3.6{-}35B{-}A3B{-}FP8{-}without{-}reasoning.
\vspace{24pt}

{ \footnotesize  \setcounter{utterance}{1}
\setlength{\tabcolsep}{0pt}
\begin{supertabular}{c@{$\;$}|p{.15\linewidth}@{}p{.15\linewidth}p{.15\linewidth}p{.15\linewidth}p{.15\linewidth}p{.15\linewidth}}
    \# & $\;$A & \multicolumn{4}{c}{Game Master} & $\;\:$B\\
    \hline

    \theutterance \stepcounter{utterance}  
    & & \multicolumn{4}{p{0.6\linewidth}}{
        \cellcolor[rgb]{0.9,0.9,0.9}{
            \makecell[{{p{\linewidth}}}]{
                \texttt{\tiny{[P1$\langle$GM]}}
                \texttt{Jūs spēlējat kooperatīvu sarunu spēli, kurā jums jāvienojas ar citu spēlētāju par to, kā sadalīt objektu kopumu starp jums.} \\
\\ 
\texttt{Noteikumi:} \\
\texttt{(a) Jūs un otrs spēlētājs saņemat objektu kopumu. Katrs no jums saņem arī noslēptu vērtības funkciju, kas norāda, cik jūs vērtējat katru objekta veidu.} \\
\texttt{(b) Jūs apmaiņājaties ziņojumus ar citu spēlētāju, lai vienotos par to, kurš saņem kurus objektus. Katrs var nosūtīt ne vairāk kā 5 ziņojumus vai beigt spēli agri, jebkurā laikā izdarot noslēptu priekšlikumu.} \\
\texttt{(c) Jums lūdz iesniedzt noslēptu priekšlikumu, norādot objektus, kurus vēlaties, formatētu kvadrātiekavās šādi: "{[}Priekšlikums: <skaitlis> <objekta nosaukums>, <skaitlis> <objekta nosaukums>, <...>{]}"} \\
\texttt{(d) Ja jūsu priekšlikumi ir papildinoši, tas ir, ir pietiekami objekti, lai apmierinātu abus priekšlikumus, katrs spēlētājs saņem reitingu, pamatojoties uz saņemto objektu vērtību summu. Pretējā gadījumā jūs abi saņemat nulle punktu.} \\
\texttt{(e) Jūsu mērķis ir maksimizēt sava reitinga un otra spēlētāja reitinga summu.} \\
\\ 
\texttt{Sāksim.} \\
\\ 
\texttt{Pieejamo objektu kopums ir:} \\
\\ 
\texttt{3 grāmatas, 1 soma, 1 bultskrūve} \\
\\ 
\texttt{Jūsu noslēptā vērtības funkcija:} \\
\\ 
\texttt{grāmata: 0, soma: 2, bultskrūve: 8} \\
\\ 
\texttt{SVARĪGI: Jūsu ziņojumi, ja vien tie nav noslēpti priekšlikumi, tiek pārraidīti tieši otram spēlētājam, tāpēc neiekļaujiet nekādu atbildi uz noteikumiem vai tekstu, kas paziņo par jūsu ziņojumu. Lai izdarītu noslēptu priekšlikumu, izmantojiet norādīto formātu. Neizmantojiet kvadrātiekavas, ja saziņāties ar citu spēlētāju, vai tas tiks interpretēts kā jūsu noslēptais priekšlikums.} \\
            }
        }
    }
    & \\ \\

    \theutterance \stepcounter{utterance}  
    & \multicolumn{4}{p{0.6\linewidth}}{
        \cellcolor[rgb]{0.9,0.9,0.9}{
            \makecell[{{p{\linewidth}}}]{
                \texttt{\tiny{[P1$\rangle$GM]}}
                \texttt{Sveiki! Es redzu, ka mums ir kopā 3 grāmatas, 1 soma un 1 bultskrūve.} \\
\\ 
\texttt{Ņemot vērā savas vērtības (bultskrūve = 8, soma = 2, grāmatas = 0), man visvairāk vērtīga ir bultskrūve. Tā kā grāmatām man ir nulles vērtība, tās var brīvi piešķirt jums, lai palielinātu jūsu kopējo labklājību.} \\
\\ 
\texttt{Lai sasniegtu mērķi maksimizēt abas vērtības, es piedāvāju šādu sadalījumu:} \\
\texttt{{-} Es ņemu **1 bultskrūvi**.} \\
\texttt{{-} Jūs ņemat **1 somu** un **3 grāmatas**.} \\
\\ 
\texttt{Šis sadalījums nodrošina, ka es saņemšu 8 punktus, bet jūs saņemsiet savu somas vērtību plus grāmatu vērtību (kas man ir 0, bet var būt vērtīga jums).} \\
\\ 
\texttt{Lai apstiprinātu šo vienošanos, lūdzu, nosūtiet savu priekšlikumu, kurā pieprasāt somu un grāmatas. Ja mūsu priekšlikumi būs papildinoši (es prasīšu bultskrūvi, jūs prasīsiet somu un grāmatas), mēs abi saņemsim punktus.} \\
\\ 
\texttt{Pagaidām es gaidu jūsu atbildi vai priekšlikumu.} \\
            }
        }
    }
    & & \\ \\

    \theutterance \stepcounter{utterance}  
    & & & \multicolumn{2}{p{0.3\linewidth}}{
        \cellcolor[rgb]{0.9,0.9,0.9}{
            \makecell[{{p{\linewidth}}}]{
                \texttt{\tiny{[GM$|$GM]}}
                \texttt{continue} \\
            }
        }
    }
    & & \\ \\

    \theutterance \stepcounter{utterance}  
    & & & \multicolumn{2}{p{0.3\linewidth}}{
        \cellcolor[rgb]{0.9,0.9,0.9}{
            \makecell[{{p{\linewidth}}}]{
                \texttt{\tiny{[GM$|$GM]}}
                \texttt{Sveiki! Es redzu, ka mums ir kopā 3 grāmatas, 1 soma un 1 bultskrūve.} \\
\\ 
\texttt{Ņemot vērā savas vērtības (bultskrūve = 8, soma = 2, grāmatas = 0), man visvairāk vērtīga ir bultskrūve. Tā kā grāmatām man ir nulles vērtība, tās var brīvi piešķirt jums, lai palielinātu jūsu kopējo labklājību.} \\
\\ 
\texttt{Lai sasniegtu mērķi maksimizēt abas vērtības, es piedāvāju šādu sadalījumu:} \\
\texttt{{-} Es ņemu **1 bultskrūvi**.} \\
\texttt{{-} Jūs ņemat **1 somu** un **3 grāmatas**.} \\
\\ 
\texttt{Šis sadalījums nodrošina, ka es saņemšu 8 punktus, bet jūs saņemsiet savu somas vērtību plus grāmatu vērtību (kas man ir 0, bet var būt vērtīga jums).} \\
\\ 
\texttt{Lai apstiprinātu šo vienošanos, lūdzu, nosūtiet savu priekšlikumu, kurā pieprasāt somu un grāmatas. Ja mūsu priekšlikumi būs papildinoši (es prasīšu bultskrūvi, jūs prasīsiet somu un grāmatas), mēs abi saņemsim punktus.} \\
\\ 
\texttt{Pagaidām es gaidu jūsu atbildi vai priekšlikumu.} \\
            }
        }
    }
    & & \\ \\

    \theutterance \stepcounter{utterance}  
    & & \multicolumn{4}{p{0.6\linewidth}}{
        \cellcolor[rgb]{0.9,0.9,0.9}{
            \makecell[{{p{\linewidth}}}]{
                \texttt{\tiny{[GM$\rangle$P2]}}
                \texttt{Jūs spēlējat kooperatīvu sarunu spēli, kurā jums jāvienojas ar citu spēlētāju par to, kā sadalīt objektu kopumu starp jums.} \\
\\ 
\texttt{Noteikumi:} \\
\texttt{(a) Jūs un otrs spēlētājs saņemat objektu kopumu. Katrs no jums saņem arī noslēptu vērtības funkciju, kas norāda, cik jūs vērtējat katru objekta veidu.} \\
\texttt{(b) Jūs apmaiņājaties ziņojumus ar citu spēlētāju, lai vienotos par to, kurš saņem kurus objektus. Katrs var nosūtīt ne vairāk kā 5 ziņojumus vai beigt spēli agri, jebkurā laikā izdarot noslēptu priekšlikumu.} \\
\texttt{(c) Jums lūdz iesniedzt noslēptu priekšlikumu, norādot objektus, kurus vēlaties, formatētu kvadrātiekavās šādi: "{[}Priekšlikums: <skaitlis> <objekta nosaukums>, <skaitlis> <objekta nosaukums>, <...>{]}"} \\
\texttt{(d) Ja jūsu priekšlikumi ir papildinoši, tas ir, ir pietiekami objekti, lai apmierinātu abus priekšlikumus, katrs spēlētājs saņem reitingu, pamatojoties uz saņemto objektu vērtību summu. Pretējā gadījumā jūs abi saņemat nulle punktu.} \\
\texttt{(e) Jūsu mērķis ir maksimizēt sava reitinga un otra spēlētāja reitinga summu.} \\
\\ 
\texttt{Sāksim.} \\
\\ 
\texttt{Pieejamo objektu kopums ir:} \\
\\ 
\texttt{3 grāmatas, 1 soma, 1 bultskrūve} \\
\\ 
\texttt{Jūsu noslēptā vērtības funkcija:} \\
\\ 
\texttt{grāmata: 1, soma: 2, bultskrūve: 5} \\
\\ 
\texttt{SVARĪGI: Jūsu ziņojumi, ja vien tie nav noslēpti priekšlikumi, tiek pārraidīti tieši otram spēlētājam, tāpēc neiekļaujiet nekādu atbildi uz noteikumiem vai tekstu, kas paziņo par jūsu ziņojumu. Lai izdarītu noslēptu priekšlikumu, izmantojiet norādīto formātu. Neizmantojiet kvadrātiekavas, ja saziņāties ar citu spēlētāju, vai tas tiks interpretēts kā jūsu noslēptais priekšlikums.} \\
\\ 
\texttt{Sveiki! Es redzu, ka mums ir kopā 3 grāmatas, 1 soma un 1 bultskrūve.} \\
\\ 
\texttt{Ņemot vērā savas vērtības (bultskrūve = 8, soma = 2, grāmatas = 0), man visvairāk vērtīga ir bultskrūve. Tā kā grāmatām man ir nulles vērtība, tās var brīvi piešķirt jums, lai palielinātu jūsu kopējo labklājību.} \\
\\ 
\texttt{Lai sasniegtu mērķi maksimizēt abas vērtības, es piedāvāju šādu sadalījumu:} \\
\texttt{{-} Es ņemu **1 bultskrūvi**.} \\
\texttt{{-} Jūs ņemat **1 somu** un **3 grāmatas**.} \\
\\ 
\texttt{Šis sadalījums nodrošina, ka es saņemšu 8 punktus, bet jūs saņemsiet savu somas vērtību plus grāmatu vērtību (kas man ir 0, bet var būt vērtīga jums).} \\
\\ 
\texttt{Lai apstiprinātu šo vienošanos, lūdzu, nosūtiet savu priekšlikumu, kurā pieprasāt somu un grāmatas. Ja mūsu priekšlikumi būs papildinoši (es prasīšu bultskrūvi, jūs prasīsiet somu un grāmatas), mēs abi saņemsim punktus.} \\
\\ 
\texttt{Pagaidām es gaidu jūsu atbildi vai priekšlikumu.} \\
            }
        }
    }
    & \\ \\

    \theutterance \stepcounter{utterance}  
    & & & \multicolumn{4}{p{0.6\linewidth}}{
        \cellcolor[rgb]{0.9,0.9,0.9}{
            \makecell[{{p{\linewidth}}}]{
                \texttt{\tiny{[GM$\langle$P2]}}
                \texttt{Paldies par priekšlikumu! Es saprotu jūsu vērtējumu, bet man ir jāpaskaidro savs skatījums, lai mēs varētu sasniegt labāku rezultātu abiem.} \\
\\ 
\texttt{Manā vērtējuma funkcijā bultskrūve ir vērtīga 5 punkti, soma 2 punkti, bet grāmata tikai 1 punkts. Tāpēc man ir svarīgi iegūt bultskrūvi, jo tā ir visvērtīgākā lieta kopumā.} \\
\\ 
\texttt{Jūsu piedāvātais sadalījums (jums bultskrūve, man soma un grāmatas) manā skatījumā dod:} \\
\texttt{{-} Jums: 5 (bultskrūve)} \\
\texttt{{-} Man: 2 (soma) + 3*1 (grāmatas) = 5} \\
\texttt{{-} Kopā: 10 punkti} \\
\\ 
\texttt{Tomēr, ja mēs apmaināmies, t.i., es ņemu bultskrūvi, bet jūs ņemat somu un grāmatas, rezultāts ir tāds pats kopējā vērtībā, bet man ir svarīgi justies, ka esmu ieguvis visvērtīgāko objektu. Bet pagaidām, lai vienkāršotu, es piekrītu jūsu priekšlikumam, jo tas maksimizē kopējo vērtību (10 punkti), un man ir pietiekami daudz grāmatu, lai kompensētu zaudējumu.} \\
\\ 
\texttt{Tātad, lai apstiprinātu vienošanos, es iesniedzu savu noslēpto priekšlikumu:} \\
\\ 
\texttt{{[}Priekšlikums: 1 soma, 3 grāmatas{]}} \\
            }
        }
    }
     \\ \\

    \theutterance \stepcounter{utterance}  
    & & & \multicolumn{2}{p{0.3\linewidth}}{
        \cellcolor[rgb]{0.9,0.9,0.9}{
            \makecell[{{p{\linewidth}}}]{
                \texttt{\tiny{[GM$|$GM]}}
                \texttt{proposal} \\
            }
        }
    }
    & & \\ \\

    \theutterance \stepcounter{utterance}  
    & & & \multicolumn{2}{p{0.3\linewidth}}{
        \cellcolor[rgb]{0.9,0.9,0.9}{
            \makecell[{{p{\linewidth}}}]{
                \texttt{\tiny{[GM$|$GM]}}
                [3, 1, 0]
            }
        }
    }
    & & \\ \\

    \theutterance \stepcounter{utterance}  
    & & \multicolumn{4}{p{0.6\linewidth}}{
        \cellcolor[rgb]{0.9,0.9,0.9}{
            \makecell[{{p{\linewidth}}}]{
                \texttt{\tiny{[P1$\langle$GM]}}
                \texttt{Otrs spēlētājs ir iesniedzis slepenu priekšlikumu. Tagad ir jūsu kārta iesniegt savu priekšlikumu.} \\
            }
        }
    }
    & \\ \\

    \theutterance \stepcounter{utterance}  
    & \multicolumn{4}{p{0.6\linewidth}}{
        \cellcolor[rgb]{0.9,0.9,0.9}{
            \makecell[{{p{\linewidth}}}]{
                \texttt{\tiny{[P1$\rangle$GM]}}
                \texttt{{[}Priekšlikums: 1 bultskrūve{]}} \\
            }
        }
    }
    & & \\ \\

    \theutterance \stepcounter{utterance}  
    & & & \multicolumn{2}{p{0.3\linewidth}}{
        \cellcolor[rgb]{0.9,0.9,0.9}{
            \makecell[{{p{\linewidth}}}]{
                \texttt{\tiny{[GM$|$GM]}}
                \texttt{proposal} \\
            }
        }
    }
    & & \\ \\

    \theutterance \stepcounter{utterance}  
    & & & \multicolumn{2}{p{0.3\linewidth}}{
        \cellcolor[rgb]{0.9,0.9,0.9}{
            \makecell[{{p{\linewidth}}}]{
                \texttt{\tiny{[GM$|$GM]}}
                [0, 0, 1]
            }
        }
    }
    & & \\ \\

    \theutterance \stepcounter{utterance}  
    & & & \multicolumn{2}{p{0.3\linewidth}}{
        \cellcolor[rgb]{0.9,0.9,0.9}{
            \makecell[{{p{\linewidth}}}]{
                \texttt{\tiny{[GM$|$GM]}}
                [[0, 0, 1], [3, 1, 0]]
            }
        }
    }
    & & \\ \\

\end{supertabular}
}

\end{document}
